\subsection{Principal Ideal Domains and Unique Factorization}

\begin{definition} \label{def:principal-ideal-domain-and-ufd} \hlt{Principal Ideal Domains and Unique Factorization Domains}\\
Let $R$ be an integral domain.
\begin{enumerate}[label=(\roman*)]
\setlength{\itemsep}{0pt}
\item The ring $R$ is a \hlt{principal ideal domain (PID)} iff $(\forall\,\mathcal{I}\trianglelefteq R)(\exists\,a\in R)\,\mathcal{I}=\langle a\rangle$.
\item The ring $R$ is a \hlt{unique factorization domain (UFD)} iff both conditions hold:
\begin{enumerate}[label=(\alph*)]
\setlength{\itemsep}{0pt}
\item Every $a\in R\setminus(R^\times\cup\{0_R\})$ admits a factorization $a=\prod\limits_{i=1}^{m}p_i$ with $m\in\N$ and irreducible $p_1,\ldots,p_m\in R$.
\item If $a=\prod\limits_{i=1}^{m}p_i=\prod\limits_{j=1}^{n}q_j$ are two such factorizations, then $m=n$ and $(\exists\,\sigma\in\mathfrak{S}_m)(\forall\,i\in\{1,\ldots,m\})\,p_i\sim q_{\sigma(i)}$.
\end{enumerate}
\end{enumerate}
\end{definition}

\begin{theorem} \label{thm:integers-and-polynomial-ring-are-pids} \hlt{$\Z$ and $\mathbb{F}[x]$ Are Principal Ideal Domains}\\
Let $\mathbb{F}$ be a field.
\begin{enumerate}[label=(\roman*)]
\setlength{\itemsep}{0pt}
\item Every ideal of $\Z$ is principal.
\item If $\{0\}\neq\mathcal{I}\trianglelefteq\Z$, then $\mathcal{I}$ has the unique positive generator $d=\min(\mathcal{I}\cap\N)$.
\item Every ideal of $\mathbb{F}[x]$ is principal.
\item If $\{0\}\neq\mathcal{I}\trianglelefteq\mathbb{F}[x]$, then its unique monic generator is the unique monic element of $\mathcal{I}$ having minimal degree.
\end{enumerate}
\begin{proof}
\begin{enumerate}[label=(\roman*)]
\setlength{\itemsep}{0pt}
\item The zero ideal is $\langle0\rangle$. Let $\{0\}\neq\mathcal{I}\trianglelefteq\Z$. Then $\mathcal{I}\cap\N\neq\emptyset$, because $\mathcal{I}$ contains the positive one of $a,-a$ whenever $0\neq a\in\mathcal{I}$.\\
Choose $d=\min(\mathcal{I}\cap\N)$ by Assumption \ref{assump:well-ordering-positive-integers}. For $a\in\mathcal{I}$, choose $q,r\in\Z$ with $a=qd+r$ and $0\leq r<d$.\\
Since $r=a-qd\in\mathcal{I}$, $r=0$ by minimality of $d$. Thus $\mathcal{I}\subseteq\langle d\rangle$. The reverse inclusion holds by $d\in\mathcal{I}$.
\item Let $e$ be a positive generator of $\mathcal{I}$. Then $d\leq e$ because $e\in\mathcal{I}\cap\N$.\\
Since $d\in\langle e\rangle$, choose $k\in\Z$ with $d=ke$. Then $k\in\N$ by $d,e>0$, so $e\leq d$. Hence $e=d$.
\item The zero ideal is $\langle0\rangle$. Let $\{0\}\neq\mathcal{I}\trianglelefteq\mathbb{F}[x]$.\\
A monic element of $\mathcal{I}$ exists by scaling a nonzero element by the inverse of its leading coefficient.\\
Choose a monic $m\in\mathcal{I}$ of minimal degree by Assumption \ref{assump:well-ordering-positive-integers} for the shifted degrees.\\
For $p\in\mathcal{I}$, write $p=qm+r$, where $r=0$ or $\deg r<\deg m$, by Theorem \ref{thm:polynomial-division-algorithm}(i).\\
Since $r=p-qm\in\mathcal{I}$, a nonzero $r$ could be scaled to a monic element of smaller degree.\\
Hence $r=0$ by minimality of $m$, so $\mathcal{I}\subseteq\langle m\rangle$; the reverse inclusion holds by $m\in\mathcal{I}$.
\item Let $m'$ be another monic element of $\mathcal{I}$ with $\deg m'=\deg m$.\\
Then $m-m'\in\mathcal{I}$ is zero or has degree less than $\deg m$, so $m=m'$ by minimality of $m$.\\
If $g$ is a monic generator of $\mathcal{I}$, then $\deg m\leq\deg g$ by minimality of $m$.\\
Choose $h\in\mathbb{F}[x]$ with $m=gh$. Then $\deg m=\deg g+\deg h$, so $\deg g\leq\deg m$. Thus $g=m$ by the preceding uniqueness.
\end{enumerate}
\hlt{Structural Commentary}: Minimal representatives and division remainders identify ideal generators.\\
\hlt{Key Insight}: A nonzero smaller remainder contradicts the selected representative.
\end{proof}
\end{theorem}

\begin{corollary} \label{cor:polynomial-generated-ideal-gcd} \hlt{GCDs Generate Polynomial Ideals}\\
Let $\mathbb{F}$ be a field, let $n\in\N$, and let $0\neq p_1,\ldots,p_n\in\mathbb{F}[x]$. Set $\mathcal{I}\defeq\langle p_1,\ldots,p_n\rangle$.
\begin{enumerate}[label=(\roman*)]
\setlength{\itemsep}{0pt}
\item A monic $m\in\mathbb{F}[x]$ with $\mathcal{I}=\langle m\rangle$ exists.
\item The monic generator $m$ in (i) is unique.
\item $(\forall\,i\in\{1,\ldots,n\})\,m\mid p_i$.
\item If $q\in\mathbb{F}[x]$ satisfies $(\forall\,i\in\{1,\ldots,n\})\,q\mid p_i$, then $q\mid m$. Consequently, $m=\gcd(p_1,\ldots,p_n)$.
\end{enumerate}
\begin{proof}
\begin{enumerate}[label=(\roman*)]
\setlength{\itemsep}{0pt}
\item The ideal $\mathcal{I}$ is nonzero because $p_1\in\mathcal{I}$. A monic $m$ with $\mathcal{I}=\langle m\rangle$ exists by Theorem \ref{thm:integers-and-polynomial-ring-are-pids}(iv).
\item The monic generator is unique by Theorem \ref{thm:integers-and-polynomial-ring-are-pids}(iv).
\item Let $i\in\{1,\ldots,n\}$. Since $p_i\in\mathcal{I}=\langle m\rangle$, $m\mid p_i$ by Definition \ref{def:unit-associate-prime-irreducible-domain}(ii).
\item Suppose $(\forall\,i\in\{1,\ldots,n\})\,q\mid p_i$. Then $q\neq0$ by $p_1\neq0$ and $q\mid p_1$.\\
The inclusions $(\forall\,i\in\{1,\ldots,n\})\,\langle p_i\rangle\subseteq\langle q\rangle$ hold by Proposition \ref{prop:principal-ideals-divisibility-associates}(iii), so $\mathcal{I}\subseteq\langle q\rangle$.\\
Thus $m\in\langle q\rangle$, and $q\mid m$ by Definition \ref{def:unit-associate-prime-irreducible-domain}(ii). Therefore $m$ has the asserted gcd property.
\end{enumerate}
\hlt{Structural Commentary}: Principal generation constructs the gcd; containment verifies universality.\\
\hlt{Key Insight}: Divisibility reverses inclusion of the corresponding principal ideals.
\end{proof}
\end{corollary}

\begin{counterexample} \label{cex:two-variable-polynomial-ring-not-pid} \hlt{A Polynomial Ring That Is Not a PID}\\
The PID conclusion for $\mathbb{F}[x]$ in Theorem \ref{thm:integers-and-polynomial-ring-are-pids}(iii) does not extend to $\mathbb{F}[x,y]$. Let $\mathbb{F}$ be a field and set $\mathcal{I}\defeq\langle x,y\rangle\trianglelefteq\mathbb{F}[x,y]$.\\
Suppose $\mathcal{I}=\langle p\rangle$. Since $x\in\mathcal{I}$ and $x\neq0$, $p\neq0$. Choose $q,r\in\mathbb{F}[x,y]$ with $x=pq$ and $y=pr$.\\
As both products nonzero, $0=\deg_y p+\deg_y q$, $0=\deg_x p+\deg_x r$. Thus $p\in\mathbb{F}[x]\cap\mathbb{F}[y]=\mathbb{F}$. Since $p\neq0$, it is a unit, so $\mathcal{I}=\mathbb{F}[x,y]$.\\
This contradicts $1\notin\mathcal{I}$, because every element of $\mathcal{I}$ has zero constant term. Hence $\mathbb{F}[x,y]$ is not a PID.
\end{counterexample}

\begin{proposition} \label{prop:pid-ascending-chain-condition} \hlt{PIDs Satisfy the Ascending Chain Condition}\\
If $R$ is a PID and $(\mathcal{I}_j)_{j\in\N}$ is an ascending chain of ideals, then $(\exists\,N\in\N)(\forall\,j\in\N)\,[j\geq N\Rightarrow\mathcal{I}_j=\mathcal{I}_N]$.
\begin{proof}
Set $\mathcal{I}\defeq\bigcup\limits_{j\in\N}\mathcal{I}_j$. The set $\mathcal{I}$ is an ideal by Proposition \ref{prop:intersections-and-chain-unions-of-ideals}. Choose $a\in R$ with $\mathcal{I}=\langle a\rangle$ by Definition \ref{def:principal-ideal-domain-and-ufd}(i).\\
Choose $N\in\N$ with $a\in\mathcal{I}_N$. Then $\mathcal{I}=\langle a\rangle\subseteq\mathcal{I}_N\subseteq\mathcal{I}$. If $j\geq N$, then $\mathcal{I}_N\subseteq\mathcal{I}_j\subseteq\mathcal{I}=\mathcal{I}_N$, so $\mathcal{I}_j=\mathcal{I}_N$.\\
\hlt{Structural Commentary}: The union is principal, and its generator enters one finite stage.\\
\hlt{Key Insight}: That stage contains the entire union and stabilizes the chain.
\end{proof}
\end{proposition}

\begin{theorem} \label{thm:pid-prime-irreducible-bezout} \hlt{Prime, Irreducible, and Bezout Properties in a PID}\\
Let $R$ be a principal ideal domain.
\begin{enumerate}[label=(\roman*)]
\setlength{\itemsep}{0pt}
\item For $r\in R\setminus(R^\times\cup\{0_R\})$, $r$ is irreducible iff $\langle r\rangle$ is a maximal ideal of $R$.
\item Every irreducible element of $R$ is prime.
\item For $a,b\in R$, $(\forall\,h\in R)\,[(h\mid a\land h\mid b)\Rightarrow h\in R^\times]$ iff $(\exists\,x,y\in R)\,xa+yb=1_R$.
\end{enumerate}
\begin{proof}
\begin{enumerate}[label=(\roman*)]
\setlength{\itemsep}{0pt}
\item[$(i)$] ($\Rightarrow$) Suppose $r$ is irreducible. The ideal $\langle r\rangle$ is proper by Proposition \ref{prop:principal-ideals-divisibility-associates}(i).\\
Let $\langle r\rangle\subseteq\mathcal{I}\trianglelefteq R$, and choose $a\in R$ with $\mathcal{I}=\langle a\rangle$ by Definition \ref{def:principal-ideal-domain-and-ufd}(i).\\
Choose $x\in R$ with $r=xa$. Then $x$ or $a$ is a unit by Definition \ref{def:unit-associate-prime-irreducible-domain}(v).\\
If $x$ is a unit, then $\langle r\rangle=\langle a\rangle=\mathcal{I}$ by Proposition \ref{prop:principal-ideals-divisibility-associates}(ii). If $a$ is a unit, then $\mathcal{I}=R$ by Proposition \ref{prop:principal-ideals-divisibility-associates}(i). Thus $\langle r\rangle$ is maximal.
\item[$(i)$] ($\Leftarrow$) Suppose $\langle r\rangle$ is maximal and write $r=ab$. Then $\langle r\rangle\subseteq\langle a\rangle$, so $\langle a\rangle=\langle r\rangle$ or $\langle a\rangle=R$.\\
In the second case, $a$ is a unit by Proposition \ref{prop:principal-ideals-divisibility-associates}(i). In the first case, $a\sim r$ by Proposition \ref{prop:principal-ideals-divisibility-associates}(ii).\\
Choose $u\in R^\times$ with $a=ur$ by Definition \ref{def:unit-associate-prime-irreducible-domain}(iii); $r=urb$, $ub=1_R$ by Proposition \ref{prop:integral-domain-cancellation}. Hence $b$ is a unit, so $r$ is irreducible.
\item[$(ii)$] Let $r$ be irreducible and suppose $r\mid ab$. Set $\mathcal{K}\defeq\langle r,a\rangle$. The ideal $\langle r\rangle$ is maximal by (i), so $\mathcal{K}=\langle r\rangle$ or $\mathcal{K}=R$.\\
In the first case, $a\in\langle r\rangle$, so $r\mid a$ by Definition \ref{def:unit-associate-prime-irreducible-domain}(ii). In the second case, choose $x,y\in R$ with $1_R=xr+ya$.\\
Then $b=xrb+yab$. Since $r\mid xrb$ and $r\mid yab$, $r\mid b$. Hence $r$ is prime by Definition \ref{def:unit-associate-prime-irreducible-domain}(iv).
\item[$(iii)$] ($\Rightarrow$) Suppose $(\forall\,h\in R)\,[(h\mid a\land h\mid b)\Rightarrow h\in R^\times]$. Choose $d\in R$ with $\langle a,b\rangle=\langle d\rangle$ by Definition \ref{def:principal-ideal-domain-and-ufd}(i).\\
Since $a,b\in\langle d\rangle$, $d\mid a$, $d\mid b$ by Definition \ref{def:unit-associate-prime-irreducible-domain}(ii). Thus $d$ is a unit, so $\langle a,b\rangle=R$ by Proposition \ref{prop:principal-ideals-divisibility-associates}(i); $(\exists\,x,y\in R)\,xa+yb=1_R$.
\item[$(iii)$] ($\Leftarrow$) Choose $x,y\in R$ with $xa+yb=1_R$, and let $h\in R$ satisfy $h\mid a$ and $h\mid b$. Then $h\mid1_R$, so $h$ is a unit by Definition \ref{def:unit-associate-prime-irreducible-domain}(i).
\end{enumerate}
\hlt{Structural Commentary}: Maximality-to-primality meets Bezout identities from two-generator ideals.\\
\hlt{Key Insight}: Principal ideal containment converts factors and common divisors into units.
\end{proof}
\end{theorem}

\begin{theorem} \label{thm:pid-is-ufd} \hlt{Every PID Is a UFD}\\
Let $R$ be a principal ideal domain.
\begin{enumerate}[label=(\roman*)]
\setlength{\itemsep}{0pt}
\item Every $a\in R\setminus(R^\times\cup\{0_R\})$ admits a factorization $a=\prod\limits_{i=1}^{m}p_i$ with $m\in\N$ and irreducible $p_1,\ldots,p_m\in R$.
\item If $a=\prod\limits_{i=1}^{m}p_i=\prod\limits_{j=1}^{n}q_j$ are two such factorizations, then $m=n$ and $(\exists\,\sigma\in\mathfrak{S}_m)(\forall\,i\in\{1,\ldots,m\})\,p_i\sim q_{\sigma(i)}$.
\end{enumerate}
\begin{proof}
\begin{enumerate}[label=(\roman*)]
\setlength{\itemsep}{0pt}
\item Suppose $a$ has no irreducible factorization and set $s_1\defeq a$.\\
Recursively suppose $s_j$ is a nonzero nonunit without an irreducible factorization.\\
Since an irreducible element is itself a factorization, $s_j$ is not irreducible.\\
Choose nonunits $u_j,v_j\in R$ with $s_j=u_jv_j$ by Definition \ref{def:unit-associate-prime-irreducible-domain}(v).\\
At least one of $u_j,v_j$ is nonfactorable; otherwise their factorizations concatenate to one of $s_j$.\\
Let $s_{j+1}$ be a nonfactorable factor and let $t_j$ be the other factor, so $s_j=s_{j+1}t_j$. Both factors are nonzero because $s_j\neq0_R$.\\
Then $\langle s_j\rangle\subseteq\langle s_{j+1}\rangle$ by Proposition \ref{prop:principal-ideals-divisibility-associates}(iii). If equality held, then $s_{j+1}\sim s_j$ by Proposition \ref{prop:principal-ideals-divisibility-associates}(ii).\\
Choose $w\in R^\times$ with $s_{j+1}=ws_j$ by Definition \ref{def:unit-associate-prime-irreducible-domain}(iii). Now $s_{j+1}=ws_{j+1}t_j$, so $1_R=wt_j$ by Proposition \ref{prop:integral-domain-cancellation}.\\
This contradicts that $t_j$ is a nonunit. Hence $\langle s_1\rangle\subsetneq\langle s_2\rangle\subsetneq\cdots$, which is impossible by Proposition \ref{prop:pid-ascending-chain-condition}.\\
Therefore every nonzero nonunit has an irreducible factorization.
\item The proof is by induction on $m$. The element $p_1$ is prime by Theorem \ref{thm:pid-prime-irreducible-bezout}(ii).\\
Then $(\exists\,j\in\{1,\ldots,n\})\,p_1\mid q_j$ by induction on $n$ with the prime property. Reindex $q$-factors, choose $u\in R^\times$ with $q_1=up_1$.\\
Choose $v\in R$ with $uv=1_R$ by Definition \ref{def:unit-associate-prime-irreducible-domain}(i). Suppose first that $m=1$. Equality $1_R=u\prod\limits_{j=2}^{n}q_j$ holds by Proposition \ref{prop:integral-domain-cancellation}.\\
If $n>1$, then $1_R=q_2\bigl(u\prod\limits_{j=3}^{n}q_j\bigr)$, so $q_2\in R^\times$ by Definition \ref{def:unit-associate-prime-irreducible-domain}(i), a contradiction. Thus $n=1$ and $p_1\sim q_1$.\\
Now let $m>1$ and suppose the conclusion holds for factorizations with $m-1$ factors on the left.\\
If $n=1$, then $\prod\limits_{i=2}^{m}p_i=u$ by Proposition \ref{prop:integral-domain-cancellation}. Then $1_R=p_2\bigl(v\prod\limits_{i=3}^{m}p_i\bigr)$, so $p_2$ is a unit by Definition \ref{def:unit-associate-prime-irreducible-domain}(i), a contradiction.\\
Hence $n>1$. Then $\prod\limits_{i=2}^{m}p_i=(uq_2)\prod\limits_{j=3}^{n}q_j$ by Proposition \ref{prop:integral-domain-cancellation}.\\
The element $uq_2$ is a nonzero nonunit, because otherwise $q_2=v(uq_2)$ would be zero or a unit.\\
If $uq_2=bc$, then $q_2=(vb)c$. Thus $vb$ or $c$ is a unit by Definition \ref{def:unit-associate-prime-irreducible-domain}(v).\\
If $vb$ is a unit, then $b=u(vb)$ is a unit. Hence $uq_2$ is irreducible by Definition \ref{def:unit-associate-prime-irreducible-domain}(v).\\
By the induction hypothesis, $m-1=n-1$, and the tail factors match up to permutation and associates.\\
Therefore $m=n$, and the first match combined with the tail permutation defines $\sigma\in\mathfrak{S}_m$ with $(\forall\,i\in\{1,\ldots,m\})\,p_i\sim q_{\sigma(i)}$.
\end{enumerate}
\hlt{Structural Commentary}: ACC terminates factor descent; primality matches irreducible factors.\\
\hlt{Key Insight}: Strict principal-ideal growth obstructs endless factor refinement.
\end{proof}
\end{theorem}
